\documentclass[12pt,a4paper]{article}

\usepackage[utf8]{inputenc}
\usepackage[T1]{fontenc}
\usepackage[margin=2.5cm]{geometry}
\usepackage{hyperref}
\usepackage{microtype}
\usepackage{parskip}
\usepackage{lmodern}

\hypersetup{colorlinks=true, urlcolor=blue!60!black}

\pagestyle{empty}

\begin{document}

\begin{flushright}
Kundai Farai Sachikonye \\
Auf der Lichtung 19, 82178 Puchheim, Germany \\
\href{https://github.com/fullscreen-triangle}{github.com/fullscreen-triangle} \\
\texttt{kundai.sachikonye@bitspark.com} \\
(+49) 1626386344 \\[6pt]
18 May 2026
\end{flushright}

\vspace{1em}

Prof.\ Roberto Cerbino \\
Soft Matter Experiments (SoMeX) Lab \\
Faculty of Physics \\
University of Vienna \\
\texttt{roberto.cerbino@univie.ac.at} \\
Job-ID: 5570

\vspace{1.5em}

\textbf{Re: University Assistant (predoctoral) --- Soft Matter / Biophysics (Job-ID 5570)}

\vspace{1em}

Dear Prof.\ Cerbino,

I am writing to apply for the predoctoral position in the SoMeX Lab (Job-ID 5570).
My background is computational biology and first-principles physical modelling
rather than experimental soft matter physics, and I want to be direct about what
that means for this application: I have developed validated frameworks for fluid
dynamics, optical physics, and biological collective motion that sit inside the
lab's core research areas --- rheology and microscopy-based dynamics --- but I
have not operated a rheometer or built a DDM acquisition pipeline.
The case I want to make is that the theoretical and computational work I bring
is a genuine complement to the experimental programme, not a substitute for it.

\textbf{First-principles fluid dynamics and rheology.}
The SoMeX Lab's primary experimental method for characterising soft matter
is rheology: measuring how yield stress fluids, colloidal suspensions, and
biological gels flow and deform under applied stress.
The \href{https://gas-giant.vercel.app/simulate}{Gas Giant} framework I
developed derives fluid dynamics --- kinetic theory, Navier-Stokes, viscosity,
and diffusion --- from a single axiom (bounded phase space geometry) via
Kirchhoff circuit laws, with zero empirical parameters.
The key result for rheology:
viscosity is $\mu = \tau_c \times g$ (partition lag $\tau_c$ times coupling
strength $g$), where both quantities are derived from the spectral structure
of the system's phase space rather than fitted to experimental data.
Validated: viscosity 2.9\,\% MAE across 12 pure liquids; Poiseuille flow
0.9\,\% mean error; adiabatic expansion 0.1--0.3\,\% error; all at zero
empirical parameters.
The framework also provides a first-principles account of the yielding
transition: in a yield stress fluid, the yield stress $\sigma_y$ corresponds
to the partition floor at which the ordered phase of the bounded oscillatory
system breaks down irreversibly --- a phase transition in the spectral structure,
not an empirical threshold.
For the lab's programme on yield stress fluids and colloidal glasses, this
gives a predictive framework for rheological properties from molecular structure
that complements rheometric measurements: the model predicts what the rheometer
should measure before the experiment is run, and the measurement validates or
corrects the theoretical partition structure.

\textbf{Optical physics and DDM analysis.}
Differential Dynamic Microscopy extracts dynamic information from microscopy
movies by computing the image structure function $D(\mathbf{q}, \Delta t)$
across a range of wave vectors and lag times. The physical content of $D$
depends on the optical transfer function of the microscope and the light-matter
interaction in the sample.
The \href{https://github.com/fullscreen-triangle/lagrangian}{Lagrangian}
framework proves Beer-Lambert absorption, Kirchhoff circuit laws, and thermal
emission as formally identical operations (Resonant Stack theorem;
transmittance $T = 1/e = 0.3679$ in all representations). This provides a
first-principles model of the optical signal in a DDM experiment: the image
contrast that encodes particle dynamics is governed by the same partition
geometry that governs bulk absorption.
At a practical level, the \href{https://github.com/fullscreen-triangle/helicopter}{Helicopter}
platform implements multi-modal microscopy analysis combining fluorescence,
brightfield, phase-contrast, and spectral channels through sequential
constraint exclusion --- a framework directly applicable to building analysis
pipelines for the lab's spinning-disk and wide-field DDM acquisition systems.
I write instrument control and image analysis code in Python; the lab's
open-source fastDDM code base is the kind of project I would contribute to
and extend.

\textbf{Active matter and biological collective dynamics.}
The SoMeX Lab's work on cellular dynamics and bacterial collectives addresses
the same physical problem as the
\href{https://github.com/fullscreen-triangle/olduvai}{Olduvai} framework:
how do collective biological motion patterns arise from the local coupling
rules of individual constituents?
Olduvai's axiom is $\oint \mathbf{J}\cdot d\mathbf{A} = 0$ --- zero net
current through any closed surface in the biological motion field.
From this single conservation law, the framework derives: the decomposition
of biological motion into supraspinal and peripheral components (validated
against 86 nights of wearable sensor data, AUC\,=\,1.00 for Parkinson's /
ataxia / healthy classification); charge conservation rates $Q_\mathrm{motor}
= 290.9\,\mathrm{mC/s}$ and $Q_\mathrm{thought} = 133.5\,\mathrm{mC/s}$
from first principles.
The same conservation law structure applies to bacterial collective motion:
the current field $\mathbf{J}$ is the velocity field of the bacterial colony,
and the constraint $\oint \mathbf{J}\cdot d\mathbf{A} = 0$ is the topological
condition for closed-circuit collective motion patterns.
This provides a theoretical framework for connecting DDM-measured intermediate
scattering functions $f(\mathbf{q}, \Delta t)$ in bacterial suspensions to
the underlying conservation laws of the collective dynamics.

\textbf{Programming and data analysis.}
I work primarily in Python, with Rust for performance-critical components.
My framework implementations include instrument-level interfacing, large-scale
data pipeline construction (Ray distributed, $>$100\,GB datasets),
GPU-accelerated computation (WebGPU, compute shaders), and image analysis
(multi-modal microscopy fusion, CNN-based feature extraction).
MATLAB I have used less extensively but can learn quickly in the context
of DDM analysis.

\textbf{Background and research proposal.}
I hold an MSc in Computational Biology (Universität Tübingen) and an MSc
in Pharmaceutical Biotechnology (HAW Hamburg), and I completed doctoral
research in Computational Lipidomics at TUM. My degree titles are not
Physics, but the physical content of the first-principles frameworks I have
built --- bounded phase space geometry, Kirchhoff circuit laws, Banach
contraction mappings on thermodynamic manifolds --- is the core theoretical
language of soft matter physics applied in a different domain context.
I have attached a two-page research proposal outlining a programme of
work connecting the Gas Giant viscosity framework and the Lagrangian optical
physics model to DDM-measured rheological dynamics in model soft matter
systems (Carbopol yield stress gels, colloidal hard-sphere suspensions).
I hold two academic referees available on request.
I am legally resident in Germany with full work authorisation and can
relocate to Vienna for the June 2026 start date.
English is native.

\vspace{1.5em}
Yours sincerely,

\vspace{2.5em}
\textbf{Kundai Farai Sachikonye}

\end{document}
